\documentclass[10pt]{article}
\usepackage[margin=0.8in]{geometry}
\usepackage{amsmath,amssymb,array}
\usepackage[hidelinks]{hyperref}
\usepackage{xurl}
\newcommand{\Tr}{\operatorname{Tr}}
\begin{document}
\begin{center}\textbf{What selects the momentum row?}\end{center}

\paragraph{Result and premises.}
Use the native CAR, the finite sine stencil and the declared
spectral completion of $H_*$. A positive Hilbert norm and a
self-adjoint generator are retained. H2's momentum multiplier
role implies an affine extension before eliminating variables,
\begin{equation}
 H_{\rm ext}=H_*+\sum_{x,i}v_x^i C_{i,x}+H_{\partial,v},\qquad
 C_i=g_i+\nabla_j\Pi_i{}^j .
 \label{eq:fm-affine}
\end{equation}
The displayed Gauss realization is a choice until its source,
flux algebra and boundary charge are identified. Varying a term
$\sum v^ig_i$ alone imposes $g_i=0$; it does not determine a
flow sourced by moving matter. H2 fixes neither the microscopic
representative $g_i$ nor its action on a ruler operator.
% [X]: multiplier variation and the explicit source/constraint
% counterfamilies below; agrees with GAP_tensor1 and P16(40C)1,3,5.

\paragraph{A derived free current, with two algebraic checks.}
On the uniform infinite lattice (or its bulk/window symbol), set
$X_i=i\partial_{p_i}$ on smooth periodic momentum wavefunctions.
At $N=\psi=1$, let
$h=m\beta+\alpha^jd_j(p)$, $d_j=\sin(ap_j)/a$, with Clifford
matrices. The symmetrically allocated energy dipole gives
\begin{align}
 J_{E,i}&=i\left[h,\tfrac12\{X_i,h\}\right]
       =\frac1{2i}[X_i,h^2]
       =\frac12\partial_{p_i}(m^2+|d|^2)\notag\\
       &=Q_i(p)\,\mathbf1,\qquad Q_i(p)=\frac{\sin(2ap_i)}{2a}.
 \label{eq:fm-free}
\end{align}
Independently, the band-symbol current
$\tfrac12\{h,\partial_{p_i}h\}$ reduces by the Clifford
anticommutators to $d_i\partial_{p_i}d_i\,\mathbf1$, the same
answer. At constant positive $N,\psi=P$, the energy current is
$N^2P^{-4}Q_i$. Identifying a momentum density with this energy
current requires a stress/boost identity for the complete theory;
the free identity does not supply one. A finite periodic position
matrix has wrap-boundary terms and is not the $X_i$ used here.
% [X]: symmetric energy-dipole commutator; Clifford band-current identity.

\paragraph{The selectors and their limits.}
Within the reflection-odd polar-vector symbol class, the six
selectors of the recorded representative have the following scope.
\begin{center}
\begin{tabular}{p{0.30\linewidth}p{0.64\linewidth}}
Spin-scalar, packet-independent one-body source &
Follows for the free current in \eqref{eq:fm-free}; H2 permits
spin, species and composite contributions.\\[3pt]
Axial hopping &
Follows from the chosen separable $|d|^2$; H2 and completion
equivariance do not exclude diagonal hops.\\[3pt]
Support at most two steps &
A minimal-range choice; finite locality alone allows larger ranges.\\[3pt]
Doubling-even transport &
A choice of the same undemodulated flow in every sine valley,
realized by $Q_i$; not an abstract Gauss-law consequence.\\[3pt]
Unit valley slope &
Fixes the normalization of $v$ in the chosen operational chart.
$C_i\mapsto cC_i$, $v^i\mapsto v^i/c$ leaves the constraint
unchanged.\\[3pt]
Endpoint-symmetric placement &
Hermiticity gives the average within a real endpoint-only ansatz;
that ansatz and the exclusion of gradient improvements are choices.
\end{tabular}
\end{center}
Reflection oddness must be stated: without it,
$Q_i+c[1-\cos(2ap_i)]/a$ satisfies the literal six conditions,
including unit derivative at each valley, but is not a polar vector
under reflection of direction $i$. With oddness, axial support two
allows $a^{-1}(c_1\sin ap_i+c_2\sin2ap_i)$. Doubling evenness
sets $c_1=0$ and unit slope sets $c_2=1/2$, proving the stated
minimal-range uniqueness.

Dropping only that range choice leaves the explicit family
\begin{equation}
 Q_i^{(b)}(p)=\sum_{r=1}^{R}b_r\frac{\sin(2rap_i)}{2ra},
 \quad b_r\in\mathbb R,\quad\sum_rb_r=1,\qquad
 J_b[v]=\frac12\sum_i\{v^i,Q_i^{(b)}\}.
 \label{eq:fm-family}
\end{equation}
Every member is Hermitian, finite range, translation covariant,
reflection odd and doubling even, with unit slope at all valleys.
Its cubic error is $-(2/3)a^2p_i^3\sum_rb_rr^2$; the coefficients
remain free. Even a fixed uniform source leaves Hermitian gradient
terms such as $a^2\{\Delta v^i,Q_i\}/2$ and local divergence
improvements. Equivariant spectral filtering and completion
transport a chosen operator; they do not select among these
operators.

\paragraph{Conservation does not rescue the identification.}
For $P_f=\sum_pf(p)c_p^\dagger c_p$ and a retained quartic vertex,
\[
 [P_f,V]=\sum_{1234}(f_1+f_2-f_3-f_4)
 V_{12;34}c_1^\dagger c_2^\dagger c_4c_3.
\]
Crystal-momentum conservation does not annul this bracket. When
the vertex connects $k,k$ to $0,2k$ with distinct internal labels,
$2Q(k)-Q(2k)=4a^2k^3+O(a^4k^5)$ is nonzero. A translation-
equivariant completion at a translation-invariant background
nevertheless commutes with the exact lattice translation $T_i$,
and hence with a modular logarithm chosen as a spectral function of $T_i$,
$T_i=e^{-iaP_i^{\rm mod}}$. This conserved modular total momentum
need not be an additive local one-body density. A continuous
periodic finite-range matrix logarithm cannot reproduce the
one-site shift: its determinant exponential has zero winding,
whereas the shift determinant has nonzero winding.

Even an abstract preserving constraint algebra leaves this freedom.
On a finite-dimensional regulated source Hilbert space, choose
Hermitian matrices $P_A$ and $H_*$. Adjoin
$u_A,\pi_A=-i\partial_{u_A}$ and choose
$U(u)=e^{i\sum_Au_AP_A}$. Then
\[
 \widetilde H=UH_*U^\dagger,\qquad C_A=U\pi_AU^\dagger,
\qquad[C_A,C_B]=[C_A,\widetilde H]=0.
\]
These identities hold on the common core
$U\bigl(C_c^\infty(\mathbb R^{\#A})\otimes\mathcal H_{\rm reg}\bigr)$.
This is an auxiliary enlargement, not a native flux construction;
its existence for every seed demonstrates non-selection. A native
momentum row needs the independently specified readout/flux algebra
and its Ward action. Requiring the complete energy-current
identification, a chosen symmetric energy allocation and the
minimal free range would select the displayed free representative,
but those are additional principles. The $b_r$ family does not
collapse under H2 and equivariant completion alone.

\paragraph{SOURCES and check scope.}
The free-current and non-selection statements are reconstructed
above; no numerical result is used. The associated final2
\texttt{STATE.tex} [S] candidate is thereby established as a free
energy current, not promoted to an interacting momentum law.
\begin{enumerate}
\item Governing duty 4(6):
\path{<internal>/GOAL_alt.md};
base candidate and H2 scope:
\path{<internal>/STATE.tex}, model,
``Frozen surface'' and O2 entries.
\item Actual six-selector, current and counterfamily source:
\path{<internal>/result.md},
sections 1.1--1.5, equations (1)--(13). The reflection premise
repairs the literal six-item uniqueness statement.
\item Actual WEB40C sources, read at the stated sections:
\path{<internal>/P15_current_and_dilation.md},
sections 2--3 and 5, current distinctions and transported completion;
\path{<internal>/P16_derivation.md},
sections 1, 3--6, modular translation, coefficient/operator
distinction and declared auxiliary algebra. Its numerical and
tensor-stiffness claims are not used.
\end{enumerate}
\end{document}
